% GENERATED FILE. Edit canonical lessons, then run npm run generate:books.
% canonical-source-hash: fnv1a64:64fe51f2c077335e
% canonical-lessons: LA-C194-sententia, LA-C194-sermo, LA-C194-oratio, LA-C194-rumor, LA-C194-historia

\chapter{Opinion, Conversation, Speech, Rumour, History}
\label{ch:la-opinion-conversation-speech-rumour-history}
\hlchaptermodality{194}

\begin{chapteropening}
\textbf{By the end of this chapter:} \emph{I can say an opinion, conversation, a speech, a rumour and history.}

Complete the last lesson of chapter 194: I can say an opinion, conversation, a speech, a rumour and history.
\end{chapteropening}

\section[sententia, sententiae]{sententia, sententiae — an opinion; a sentence}

\label{lesson:LA-C194-sententia}



\begin{quote}
Before the new one: say the Latin for a bracelet, then the Latin for a necklace.
\end{quote}

\subsection*{You'll want to know: sententia, sententiae}
\textbf{sententia, sententiae} — \textquotedblleft{}an opinion; a sentence\textquotedblright{}.

\begin{cousinweb}
\textbf{sententia} is shortened from \emph{sentientia}, \textquotedblleft{}a feeling\textquotedblright{} (Lewis \& Short), built on \emph{sentiō}, \textquotedblleft{}to feel, to think\textquotedblright{}: what one thinks.
\end{cousinweb}

\begin{grammarlens}[title={What you've built}]
One more word for this chapter.
\end{grammarlens}

\subsection*{Your turn}
\begin{itemize}
\raggedright
  \item \emph{Say it:} \emph{sententia, sententiae}
  \item \emph{Say it:} \emph{sententia, sententiae}, once more
  \item \emph{Say it:} say \emph{sententia, sententiae}
  \item \emph{Recall:} say the Latin for a bracelet, then the Latin for a necklace, then say \emph{sententia, sententiae} again
\end{itemize}

\begin{tcolorbox}[breakable,colback=teal!4,colframe=teal!35!black,title={Before you move on}]
What does sententia, sententiae mean? (\textquotedblleft{}An opinion; a sentence\textquotedblright{}.) Say it once more.
\end{tcolorbox}
\section[sermō, sermōnis]{sermō, sermōnis — conversation, speech}

\label{lesson:LA-C194-sermo}



\begin{quote}
Before the new one: say the Latin for a necklace, then the Latin for an opinion; a sentence.
\end{quote}

\subsection*{You'll want to know: sermō, sermōnis}
\textbf{sermō, sermōnis} — \textquotedblleft{}conversation, speech\textquotedblright{}.

\begin{cousinweb}
\textbf{sermō} is often linked with \emph{serō}, \textquotedblleft{}to join, to string together\textquotedblright{} (Lewis \& Short): talk as words strung together. The link is likely but not certain.
\end{cousinweb}

\begin{grammarlens}[title={What you've built}]
One more word for this chapter.
\end{grammarlens}

\subsection*{Your turn}
\begin{itemize}
\raggedright
  \item \emph{Say it:} \emph{sermō, sermōnis}
  \item \emph{Say it:} \emph{sermō, sermōnis}, once more
  \item \emph{Say it:} say \emph{sermō, sermōnis}
  \item \emph{Recall:} say the Latin for a necklace, then the Latin for an opinion; a sentence, then say \emph{sermō, sermōnis} again
\end{itemize}

\begin{tcolorbox}[breakable,colback=teal!4,colframe=teal!35!black,title={Before you move on}]
What does sermō, sermōnis mean? (\textquotedblleft{}Conversation, speech\textquotedblright{}.) Say it once more.
\end{tcolorbox}
\section[ōrātiō, ōrātiōnis]{ōrātiō, ōrātiōnis — a speech}

\label{lesson:LA-C194-oratio}



\begin{quote}
Before the new one: say the Latin for an opinion; a sentence, then the Latin for conversation.
\end{quote}

\subsection*{You'll want to know: ōrātiō, ōrātiōnis}
\textbf{ōrātiō, ōrātiōnis} — \textquotedblleft{}a speech\textquotedblright{}.

\begin{cousinweb}
\textbf{ōrātiō} is built on \emph{ōrō}, \textquotedblleft{}to speak, to plead, to pray\textquotedblright{} (Lewis \& Short), with \emph{-tiō}, an ending for an act. \emph{ōrō} is often linked with \emph{ōs}, \textquotedblleft{}mouth\textquotedblright{} (EtymDB 2.1), but the link is uncertain.
\end{cousinweb}

\begin{grammarlens}[title={What you've built}]
One more word for this chapter.
\end{grammarlens}

\subsection*{Your turn}
\begin{itemize}
\raggedright
  \item \emph{Say it:} \emph{ōrātiō, ōrātiōnis}
  \item \emph{Say it:} \emph{ōrātiō, ōrātiōnis}, once more
  \item \emph{Say it:} say \emph{ōrātiō, ōrātiōnis}
  \item \emph{Recall:} say the Latin for an opinion; a sentence, then the Latin for conversation, then say \emph{ōrātiō, ōrātiōnis} again
\end{itemize}

\begin{tcolorbox}[breakable,colback=teal!4,colframe=teal!35!black,title={Before you move on}]
What does ōrātiō, ōrātiōnis mean? (\textquotedblleft{}A speech\textquotedblright{}.) Say it once more.
\end{tcolorbox}
\section[rūmor, rūmōris]{rūmor, rūmōris — a rumour}

\label{lesson:LA-C194-rumor}



\begin{quote}
Before the new one: say the Latin for conversation, then the Latin for a speech.
\end{quote}

\subsection*{You'll want to know: rūmor, rūmōris}
\textbf{rūmor, rūmōris} — \textquotedblleft{}a rumour\textquotedblright{}.

\begin{cousinweb}
\textbf{rūmor} is often traced to a Proto-Indo-European root \emph{*rewH-}, \textquotedblleft{}to roar, to shout\textquotedblright{}, but the link is not certain. Lewis \& Short put \emph{raucus}, \textquotedblleft{}hoarse\textquotedblright{}, in the same family.
\end{cousinweb}

\begin{grammarlens}[title={What you've built}]
One more word for this chapter.
\end{grammarlens}

\subsection*{Your turn}
\begin{itemize}
\raggedright
  \item \emph{Say it:} \emph{rūmor, rūmōris}
  \item \emph{Say it:} \emph{rūmor, rūmōris}, once more
  \item \emph{Say it:} say \emph{rūmor, rūmōris}
  \item \emph{Recall:} say the Latin for conversation, then the Latin for a speech, then say \emph{rūmor, rūmōris} again
\end{itemize}

\begin{tcolorbox}[breakable,colback=teal!4,colframe=teal!35!black,title={Before you move on}]
What does rūmor, rūmōris mean? (\textquotedblleft{}A rumour\textquotedblright{}.) Say it once more.
\end{tcolorbox}
\section[historia, historiae]{historia, historiae — history}

\label{lesson:LA-C194-historia}



\begin{quote}
Before the new one: say the Latin for a speech, then the Latin for a rumour.
\end{quote}

\subsection*{You'll want to know: historia, historiae}
\textbf{historia, historiae} — \textquotedblleft{}history\textquotedblright{}.

\begin{cousinweb}
\textbf{historia} was borrowed from Greek \emph{historia}, \textquotedblleft{}inquiry, a narrative\textquotedblright{} (Lewis \& Short), from \emph{histōr}, \textquotedblleft{}one who knows\textquotedblright{} (EtymDB 2.1), whose root is usually given as \emph{*weyd-}, \textquotedblleft{}to see, to know\textquotedblright{}, the root of \emph{videō}, \textquotedblleft{}to see\textquotedblright{}.
\end{cousinweb}

\begin{grammarlens}[title={What you've built}]
One more word for this chapter.
\end{grammarlens}

\subsection*{Your turn}
\begin{itemize}
\raggedright
  \item \emph{Say it:} \emph{historia, historiae}
  \item \emph{Say it:} \emph{historia, historiae}, once more
  \item \emph{Say it:} say \emph{historia, historiae}
  \item \emph{Recall:} say the Latin for a speech, then the Latin for a rumour, then say \emph{historia, historiae} again
\end{itemize}

\begin{tcolorbox}[breakable,colback=teal!4,colframe=teal!35!black,title={Before you move on}]
What does historia, historiae mean? (\textquotedblleft{}History\textquotedblright{}.) Say it once more.
\end{tcolorbox}
